% 本文件由 LaTeX 排版 Agent 根据有效 translation.json 自动生成。
% author=Ignatius of Antioch; work_id=0104; title=依纳爵致厄弗所人书

\chapter{依纳爵致厄弗所人书}

% structure: p0001 source=h1 target=omitted reason=page-title-already-rendered-by-parent

% structure: p0002 source=h2 target=section reason=relative-heading-rank; base=h2

\section{问候}

依纳爵，亦号带神者，致书于亚细亚的厄弗所教会，这教会理应是最有福的，因着天主圣父的伟大和圆满而被祝福，并在时间开始以前就被预定，要永远享有持久不变的荣耀，藉着真正受难，按圣父的旨意，并通过我们的天主耶稣基督而结合并被拣选：愿你们藉着耶稣基督和祂无玷的恩宠，满溢喜乐。\par

% structure: p0004 source=h2 target=section reason=relative-heading-rank; base=h2

\section{第一章 称赞厄弗所人}

我已熟悉你们那在至爱中、与天主密切相关的名声，这是你们因正义的习惯而获得的，按照我们对救主耶稣基督的信德和爱德。你们作为天主的追随者，藉着天主的血激励自己，已完美地完成了与你们相称的工作。因为当听到我从叙利亚被捆绑而来，为着共同的名字和希望，并相信藉着你们的祈祷，我得以被允许在罗马与野兽搏斗，以便通过殉道我才能真正成为祂的门徒——祂\BibleQuote{为我们舍弃了自己，作为献给天主的馨香祭品和祭献}（见：\BibleRef{弗 5:2}）（你们急忙来见我）。因此，我奉天主的名，藉着奥乃西摩——一个具有无法言表的爱德、且在肉身内为你们的主教的人——接待了你们全体。我藉着耶稣基督恳求你们要爱他，并都努力效法他。愿那使你们配得如此杰出主教的祂受赞美。\par

% structure: p0006 source=h2 target=section reason=relative-heading-rank; base=h2

\section{第二章 祝贺与恳求}

关于我的同仆布尔黑（他在天主方面为你们的执事，并在一切事上有福的），我请求他能多留一些时候，这既为你们的荣誉，也为你们的主教。至于克洛古，他既配得天主也配得你们，我已将他当作你们爱德的明证领受了；他在一切事上使我振作（见：\BibleRef{格前 16:18}），正如我们主耶稣基督之父也要使他振作一样（见：\BibleRef{格前 16:18}）；连同奥乃西摩、布尔黑、欧普洛、弗龙多，藉着他们，我在爱德中看到了你们全体。愿我常因你们而喜乐，若我真配得的话。因此，你们理当在一切事上荣耀耶稣基督，因为祂已荣耀了你们，好使你们藉着同心合意的服从，\BibleQuote{在相同的心思和相同的判断上完全结合，一致地就同一件事说出同一句话}（见：\BibleRef{格前 1:10}），并且，服从主教和长老团，你们就能在各方面被圣化。\par

% structure: p0008 source=h2 target=section reason=relative-heading-rank; base=h2

\section{第三章 劝勉合一}

我不命令你们，好像我是什么了不起的人物。因为我虽为（基督的）名被捆绑，但在耶稣基督内我尚未成全。因为现在我刚开始作门徒，并向你们说话，你们是我的同门徒。因为你们在信德、劝勉、忍耐和宽宏方面曾激励我，这正是我所需的。但因爱德不容我对你们沉默，我便率先劝勉你们，要大家一致按天主的旨意奔跑。因为甚至耶稣基督，我们不可分离的生命，就是圣父（显明）的旨意；同样，各处安置到地极的主教，也是出于耶稣基督的旨意。\par

% structure: p0010 source=h2 target=section reason=relative-heading-rank; base=h2

\section{第四章 继续同一主题}

因此，你们应当按你们主教的旨意一同奔跑，这事你们也确实在行。因为你们那理应受赞颂、配得天主的长老团，与主教的配合恰如琴弦与竖琴一般。因此，在你们的和谐与相爱中，耶稣基督被歌颂。你们各人成为一合唱团，好使在爱中和谐，齐声唱出天主的歌，同声藉着耶稣基督向圣父歌唱，使祂能听见你们，并从你们的作为看出你们真是祂圣子的肢体。因此，活在无可指责的合一中是有益的，这样你们就常能享受与天主的共融。\par

% structure: p0012 source=h2 target=section reason=relative-heading-rank; base=h2

\section{第五章 称扬合一}

因为若我在如此短暂的时间里，与你们的主教享有如此深厚的共融——我指的不是属人的，而是属神的共融——那么我认为你们更是有福的，你们与他如此紧密相连，如同教会与耶稣基督相连，耶稣基督与圣父相连，好使万事在合一中一致！谁也不要自欺：若有人不在祭台内，他就被剥夺了天主的食粮。因为若一两个人的祈祷有如此能力（见：\BibleRef{玛 18:19}），那么主教和整个教会的祈祷更具有何等能力！因此，那不与教会聚会的人，正以此显出了自己的骄傲，并定了自己的罪。因为经上记载：\Quote{天主拒绝骄傲的人。}因此，我们要谨慎，不要与主教对立，好使我们服从天主。\par

% structure: p0014 source=h2 target=section reason=relative-heading-rank; base=h2

\section{第六章 尊敬主教如同基督本人}

一个人越看到主教保持沉默，就越应敬畏他。因为对于家主派来管理祂家室的人，我们应当接待，如同接待那差遣他的人一样（见：\BibleRef{玛 24:25}）。因此，显然我们应看待主教如同看待主本人。实际上，奥乃西摩本人极大地称赞你们在天主内的良好秩序：你们都按真理生活，没有任何派别在你们中间立足。确实，你们不听任何人超过听讲真理的耶稣基督。\par

% structure: p0016 source=h2 target=section reason=relative-heading-rank; base=h2

\section{第七章 提防假教师}

因为有些人惯于邪恶地携带（耶稣基督）的名字，却行不配天主的事，你们必须逃避他们如同逃避野兽。他们是贪食的狗，暗中咬人，你们必须提防他们，因为他们是几乎无法医治的人。只有一位医生，祂兼具肉身和神魂；既是受造的，又是不受造的；天主存在于肉身；在死亡中的真实生命；既出于玛利亚，又出于天主；先可受苦，后不可受苦——就是我们的主耶稣基督。\par

% structure: p0018 source=h2 target=section reason=relative-heading-rank; base=h2

\section{第八章 再次称赞厄弗所人}

那么，不要让任何人欺骗你们——其实你们并未受骗，因为你们完全归属于天主。既然你们中间没有纷争困扰你们，你们确实按天主的旨意生活。我远不及你们，需要由你们那闻名天下的\Place{厄弗所}教会来圣化我。属血肉的人不能做属灵的事，属灵的人也不能做属血肉的事；正如信德不能做不信的工，不信也不能做信德的工。但你们按肉体所做的事也是属灵的，因为你们在耶稣基督内做一切事。\par

% structure: p0020 source=h2 target=section reason=relative-heading-rank; base=h2

\section{第九章 你们没有听从假教师}

然而，我听说有些带着假教义的人从那里来到你们当中，你们没有容许他们在你们中间播种，却掩耳不听，免得你们接受他们所播下的东西，如同圣父之殿的石块\BibleRef{伯前 2:5}，预备为天主圣父的建筑，并由耶稣基督的工具——就是十字架——高举上升\BibleRef{若 12:32}，以圣神为绳索，你们的信德是上升的途径，而你们的爱德是通往天主的道路。因此，你们和所有旅途同伴都是背负天主者、背负圣殿者、背负基督者、背负圣洁者，在各方面以耶稣基督的诫命为装饰。在这位内，我也欢欣，因我得以藉这封信与你们交谈同乐，因为就你们的基督徒生活而言，你们除天主外不爱任何事物。\par

% structure: p0022 source=h2 target=section reason=relative-heading-rank; base=h2

\section{第十章 劝勉祈祷、谦卑等}

并且不断为他人祈祷。因为他们有悔改的希望，可以归向天主。那么，看吧，如果他们别无他法，就让他们因你们的行为而受教。面对他们的忿怒，要温柔；面对他们的自夸，要谦卑；面对他们的亵渎，以祈祷回报；面对他们的错谬，要在信德上坚定不移\BibleRef{哥 1:23}；面对他们的残暴，要显出你们的温和。我们小心不去效法他们的行为，却要在一切真正的友善中被视为他们的弟兄；让我们寻求效法主（谁曾受过比祂更不公正的待遇，更贫穷，更被定罪呢？），好使魔鬼的苗裔不在你们中间，而使你们在耶稣基督内，无论肉体或灵魂，都保持完全的圣洁和清醒。\par

% structure: p0024 source=h2 target=section reason=relative-heading-rank; base=h2

\section{第十一章 劝勉敬畏天主等}

末期已来到我们身上。因此，让我们怀着敬畏的心，畏惧天主的忍耐，免得它导致我们的定罪。因为我们要么敬畏那将要来临的忿怒，要么顾念那现今彰显的恩宠——二者必居其一。无论如何，愿我们在基督耶稣内被找到，得享真正的生命。除祂以外，不要让任何事物吸引你们，我为祂佩戴着这些锁链，这些属灵的珍宝；愿藉着你们的祈祷，我能从其中兴起，我恳求常常有分于这祈祷，好使我在\Place{厄弗所}的基督徒的份中被找到，他们因耶稣基督的能力，始终与宗徒们同心合意。\par

% structure: p0026 source=h2 target=section reason=relative-heading-rank; base=h2

\section{第十二章 赞扬厄弗所人}

我知道我是谁，也知道我写给谁。我是一个被定罪的人，你们是蒙怜悯的对象；我处于危险之中，你们却坚定在安全中。你们是那些为天主而被剪除的人所经过的人。你们与\Person{保禄}——那位圣者、殉道者、理应最为有福者——一同被引入福音的奥秘；当我得以到达天主那里时，愿我能在他的脚前被找到；他在所有的书信中都因基督耶稣提到你们。\par

% structure: p0028 source=h2 target=section reason=relative-heading-rank; base=h2

\section{第十三章 频繁聚集崇拜天主}

那么，要留意常常聚集在一起感谢天主，并显扬祂的赞美。因为当你们频繁在同一地点聚会时，\Person{撒殚}的权势就被摧毁，而他意图的毁灭也被你们信德的合一所阻止。没有比和平更珍贵的了，因着和平，天上地上的一切战争都得以终止。\par

% structure: p0030 source=h2 target=section reason=relative-heading-rank; base=h2

\section{第十四章 劝勉信德和爱德}

如果你们完全拥有那对耶稣基督的信德和爱德\BibleRef{弟前 1:14}——它们是生命的起始和终结——这些事没有一件对你们是隐藏的。因为起始是信德，终结是爱德\BibleRef{弟前 1:5}。这两者不可分离地相连，来自天主，而所有其他为圣洁生活所需的事都随着它们而来。没有人真实地宣认信德而仍然犯罪\BibleRef{若一 3:7}；拥有爱德的人也不仇恨任何人。树由果实显明\BibleRef{玛 12:33}；同样，那些自称基督徒的人必由他们的行为被认出。因为现在所要求的不是单纯的宣认，而是人要持续在信德的能力中直到终局。\par

% structure: p0032 source=h2 target=section reason=relative-heading-rank; base=h2

\section{第十五章 劝勉以沉默和言语宣认基督}

一个人沉默而成为基督徒，比谈论而并非基督徒更好。教训人是好的，如果说话的人也行动。只有一位教师，祂一说话，事就成了；甚至祂在沉默中所行的事也堪当圣父。拥有耶稣话语的人，实在能听见祂的沉默，好使他得以成全，并言行一致，且因他的沉默而被认出。没有一件事向天主隐藏，我们的秘密都临近祂。因此，我们当像有祂住在我们内的人一样做一切事，使我们成为祂的殿宇\BibleRef{格前 6:19}，祂在我们内作我们的天主——祂实在就是如此——并且要在我们面前显示自己。因此我们正当地爱祂。\par

% structure: p0034 source=h2 target=section reason=relative-heading-rank; base=h2

\section{第十六章 假教师的命运}

我的弟兄们，不要错误\BibleRef{雅 1:16}。败坏家庭的人不能继承天主的国\BibleRef{格前 6:9}。如果那些在肉体上这样做的人已遭受死亡，那以邪恶教义败坏天主的信德——耶稣基督曾为其被钉十字架——的人，更要如何呢！这样的人如此被玷污，将进入永火，凡听从他的每一个人也是如此。\par

% structure: p0036 source=h2 target=section reason=relative-heading-rank; base=h2

\section{第十七章 谨防假教义}

为此，主才允许香膏倒在祂头上（\BibleRef{若 12:7}），使祂能将不朽吹入祂的教会。你们不要被这世界之王的教义的恶臭所傅油；不要让他将你们从摆在面前的生命中掳去。我们既然领受了认识天主的恩赐，即耶稣基督，为何不全然明智？为何我们愚昧地丧亡，不承认主真正赐给我们的恩赐？\par

% structure: p0038 source=h2 target=section reason=relative-heading-rank; base=h2

\section{第十八章 十字架的光荣}

愿我的灵魂为十字架的缘故被视为无有；十字架为不信者是绊脚石（\BibleRef{格前 1:18}），但为我们却是救恩和永生。\Quote{智者在哪里？辩士在哪里？}（\BibleRef{格前 1:20}）那些被称为明智者的夸耀在哪里？因为我们的天主耶稣基督，按照天主的安排，由\Person{玛利亚}从\Person{达味}的后裔受孕，却是因圣神而成。祂出生并受洗，为藉祂的苦难洁净水。\par

% structure: p0040 source=h2 target=section reason=relative-heading-rank; base=h2

\section{第十九章 三个著名的奥秘}

\Person{玛利亚}的童贞和她的产子，以及主的死亡，这三件著名的奥秘，由天主在静默中完成，原是对这世界之主隐藏的。那么，祂是如何向世界显示的呢？一颗星在天上闪耀，超越众星，其光不可言喻，它的新奇令人惊异。其余众星连同日月都向这颗星形成合唱团，它的光辉远超它们之上。人们感到骚动，不知这新的景象从何而来，它与天上的其他一切如此不同。于是，各种巫术被摧毁，一切邪恶的枷锁消失；无知被除去，旧的王国被废除，天主亲自以人形显现，为更新永生。那由天主所预备的，如今开始应验。从此，万物陷入动荡，因为祂默想废除死亡。\par

% structure: p0042 source=h2 target=section reason=relative-heading-rank; base=h2

\section{第二十章 另一封信的许诺}

若耶稣基督因你们的祈祷而恩准，并合乎祂的旨意，我将在写给你们的第二篇小作中，进一步向你们阐明我已开始谈论的关于新人耶稣基督在祂的信德、爱德、苦难和复活中的安排。尤其若主使我得知你们各人藉着恩宠，在同一信德内，并在按肉身是\Person{达味}后裔、为人子亦为天主的耶稣基督内，共同聚集，以不分的心服从主教和长老团，擘开同一个饼，那就是不朽的良药、阻止我们死亡的解毒剂，并使我们因耶稣基督而永远活着。\par

% structure: p0044 source=h2 target=section reason=relative-heading-rank; base=h2

\section{第二十一章 结语}

愿我的灵魂为你们的灵魂和那些因天主的荣耀而被你们送到\Place{斯米纳}的人的灵魂；我也从那里写信给你们，感谢主，并爱\Person{波利卡普}如同爱你们一样。请记念我，如同耶稣基督也记念了你们。请为\Place{叙利亚}的教会祈祷，我由此被捆锁带到\Place{罗马}，我是那里末后的信徒，正如我被认为配得被选来彰显天主的荣耀。愿你们在天主圣父和我们的共同希望耶稣基督内平安。\par

\SourceRecord{Translated by Alexander Roberts and James Donaldson.；Ante-Nicene Fathers；Vol. 1.；Edited by Alexander Roberts, James Donaldson, and A. Cleveland Coxe.；Buffalo, NY: Christian Literature Publishing Co.,；1885.；<http://www.newadvent.org/fathers/0104.htm>.}
